\label{sec:literature_review}
\subsection{Non-imaging classification and feature reduction}
Clinical, biochemical, and bioimpedance measurements have been investigated as non-imaging predictors of existing gallstone status. \textcite{esen2024early} introduced Gallstone-1, comprising 319 participants (161 cases and 158 controls) and 38 predictors, now available through UCI \parencite{uci2025gallstone}. Among nine conventional classifiers, gradient boosting achieved the highest reported accuracy (85.42\%). Vitamin D, C-reactive protein (CRP), total body water, and lean mass emerged as important predictors. These findings support the predictive relevance of heterogeneous measurements and motivate nonlinear modelling and feature reduction; they do not establish prospective disease prediction.

Subsequent optimisation studies show that dimensionality reduction offers a more consistent benefit than increased accuracy. \textcite{sarker2025gallstone} compared random forest with and without Sand Cat Swarm Optimization, under holdout and ten-fold cross-validation. Using 38 features, holdout accuracy, F1, precision, and recall were 81.25\%, 79.07\%, 85.00\%, and 73.91\%, respectively; mean cross-validation accuracy was 78.42\%. Optimisation retained 13 predictors, with 79.17\% holdout accuracy and 78.32\% cross-validation accuracy. Thus, approximately two thirds of the measurements were removed while maintaining similar performance. SHAP identified CRP, vitamin D, and aspartate aminotransferase as influential, while DiCE supplied individual counterfactual explanations.

\textcite{shrestha2026metaheuristic} combined Rotation Forest with Bald Eagle Search. The full-feature baseline achieved 78.00\% accuracy, 77.35\% F1, 79.33\% precision, 75.79\% recall, and approximately 0.867 AUC. Retaining 17 predictors reduced accuracy to 75.79\% and AUC to approximately 0.860. Near-perfect training performance contrasted with substantially lower test scores, indicating limited generalisation despite model complexity. SHAP and LIME highlighted CRP, vitamin D, obesity-related measurements, haemoglobin, and body mass. Feature reduction therefore simplified measurement requirements without improving discrimination.

\textcite{das2026sloa} reported a modest benefit from Sea Lion Optimization with CatBoost. Five-fold mean accuracy increased from 79.58\% to 80.42\% while the feature count fell from 38 to 19; F1 increased from 79.01\% to approximately 79.94\%, precision from 80.91\% to 81.97\%, and recall from 77.36\% to 77.76\%. Optimised fold accuracy ranged from approximately 71.88\% to 87.50\%, demonstrating sensitivity to partitioning. SHAP, LIME, and DiCE again emphasised CRP and vitamin D. Collectively, these studies favour assessing measurement burden, stability, and generalisation alongside headline accuracy.

\subsection{Consensus selection, ensembles, and probability assessment}
\textcite{tamer2026novel} combined four statistical selectors (minimum redundancy maximum relevance, chi-square, ANOVA, and Kruskal--Wallis) with four metaheuristics (Grey Wolf, Whale, Harris Hawks, and Particle Swarm optimisation). Agreement among at least seven selectors retained 11 predictors. Across nine classifiers, random forest achieved the highest reported accuracy (90.62\%); CRP, haemoglobin, vitamin D, and hyperlipidaemia remained in the subset. Consensus selection may identify complementary information and reduce redundancy, but its stability requires direct assessment. Selection informed by observations outside the training partition would yield optimistic estimates; repeated or nested validation is therefore necessary to evaluate the procedure reliably.

Hyperparameter optimisation provides a distinct route to improving performance. \textcite{rahman2026hbb} combined Bees Algorithm global search with Bayesian optimisation for XGBoost. HBB-Opt achieved 87.50\% accuracy, 87.49\% F1, 0.9258 ROC--AUC, 90.62\% sensitivity, and 90.00\% negative predictive value. A Brier score of 0.098 additionally assessed probabilistic accuracy, although this score alone does not establish calibration. SHAP highlighted CRP, vitamin D, obesity percentage, visceral fat area, and extracellular-fluid/total-body-water ratio. This work broadens evaluation beyond class labels to the quality of predicted probabilities.

\textcite{indirani2026ensemble} investigated stacking and soft voting across random forest, XGBoost, logistic regression, gradient boosting, Extra Trees, LightGBM, and support vector machines, incorporating SHAP selection, calibration, and threshold adjustment. Random-forest/XGBoost stacking with a logistic meta-learner achieved 85.94\% accuracy and 0.9307 AUC. Soft-voting accuracy increased from 79.69\% at the default threshold to 87.50\% after validation-based adjustment, with approximately 94\% positive-class recall. SHAP-based selection achieved approximately 0.9238 AUC and the same adjusted accuracy. These results distinguish probability ranking from operating-point decisions: strong AUC or accuracy does not necessarily imply adequate sensitivity. Threshold selection must remain within development data, and calibration requires separate evaluation.

\subsection{Synthesis and study positioning}
Across these studies, tree-based ensembles capture useful predictive patterns, but additional optimisation does not consistently improve generalisation. CRP, vitamin D, and body-composition measurements recur across attribution and selection analyses, supporting their predictive relevance within this cohort. Such recurrence is not independent biological validation: most studies reuse the same 319 observations, and neither attribution nor counterfactual explanation establishes causality. Near-perfect training results in random forest, Rotation Forest, and CatBoost further highlight overfitting risk when predictors are numerous relative to sample size.

Differences in partitions, selection procedures, search budgets, and thresholds preclude ranking methods directly from published scores. Larger independent cohorts remain necessary to establish transferability, calibration, and clinical utility. The present study consequently compares an interpretable logistic-regression baseline with nonlinear and stacked alternatives under a shared, training-fitted pipeline and leave-one-patient-out (LOPO) evaluation. Accuracy, AUC, sensitivity, specificity, F1, and false-negative behaviour characterise complementary aspects of performance. LOPO supplies one held-out prediction per record while retaining almost the entire small cohort for training; it provides internal validation rather than resolving external generalisation or model-selection bias.
